\section{What Stays: The Business Rules}\label{sec:keeps}

\begin{plainwords}
Everything in this section comes straight from the v11 prototype and is treated as a
non-negotiable requirement. The new system must behave the same way -- same gates, same
approvals, same queues -- so that staff recognise it immediately. If a rule is ever
changed, it is a deliberate business decision, never an accident of the rebuild.
\end{plainwords}

\subsection{Customer master and order snapshots}

\begin{tabularx}{\textwidth}{@{}X X@{}}
\toprule
\textbf{Rule} & \textbf{Why it matters} \\
\midrule
Customer master stores name, contact, email, phone, billing address, default dispatch
  address, default dispatch method, account reference and notes &
  No re-typing for repeat customers \\
Selecting a customer copies their defaults onto the job as a \textbf{snapshot} &
  Every order is frozen as it was taken \\
Job contact/address/method are snapshots, not live links &
  Old orders never change when the customer moves house \\
``Deliberately update customer master'' is a separate, confirmed, audited action &
  Routine order editing cannot corrupt master data \\
\bottomrule
\end{tabularx}

\textbf{Database consequence:} \path{customers} table is separate from
\path{job_contact_snapshot} and \path{job_dispatch_snapshot}. The update-master
action is its own API call that always writes an audit entry.

\subsection{Job structure and product grid}

\begin{itemize}
  \item Job fields: Customer PO, Print/Job Name, Order Date, Order Type (Bulk, POD,
        Repeat, Sample), Priority (1 = highest), assigned staff, process date,
        dispatch date/time.
  \item Size grid XS--5XL per product/colour line, stored as a structured table.
  \item SKU = Manufacturer/Master Product Code; Supplier SKU kept as a separate field.
  \item \textbf{Outstanding quantity = ordered $-$ received}, computed from stock events,
        negative values allowed (over-delivery is real and must show).
  \item Search by job number, customer, PO, print name, master SKU, supplier SKU.
  \item Queue order: dispatch deadline, then priority, then process date.
\end{itemize}

\subsection{Commercial visibility (money is compartmentalised)}

\begin{tabular}{@{}p{3.9cm} p{3.4cm} p{4.4cm} p{3.6cm}@{}}
\toprule
\textbf{Field} & \textbf{Enter} & \textbf{View} & \textbf{Never sees} \\
\midrule
Buying costs & Admin, Operations & Admin, Operations; Director in reports & Office,
  Dispatch, Packing, Department Operator \\
Selling price / tax & authorised commercial roles only & same set & everyone else \\
\bottomrule
\end{tabular}

These are \textbf{field-level server-side permissions}: the API removes restricted fields
from every response. Hiding columns in the interface is not enough.

\subsection{Artwork approval}

\begin{itemize}
  \item Per job: proof reference, version, approval status, approved by/date, Pantone
        notes, uploads, links, positions.
  \item Only Admin and Operations approve or revise approved artwork. Everyone else is
        read-only.
  \item \textbf{Revising approved artwork resets it to Draft} and writes an audit entry
        with before/after -- the revised work must be approved again.
  \item Operators see production references and approval state, but cannot change them.
\end{itemize}

\textbf{State machine:}
\begin{center}
  \texttt{Draft $\rightarrow$ Awaiting Approval $\rightarrow$ Approved \textbar{} Rejected}
  \qquad
  \texttt{Approved $\xrightarrow{revision}$ Draft (new version, audited)}
\end{center}

\subsection{Embroidery swatching -- the most complex workflow}

\begin{plainwords}
Embroidery jobs need an approved sample (swatch) before full production starts. The
embroiderers make a sample, it gets approved or rejected with a reason, rejected samples
stay visible for learning, and approved samples can never be edited. If a required sample
is not yet approved, the server refuses to let the embroidery stage start, progress or
finish. Operations can waive the requirement, but only with a recorded reason.
\end{plainwords}

\begin{itemize}
  \item Swatch required by default for embroidery jobs; waiver by Admin/Ops with audit.
  \item Each attempt records: sample quantity, embroidery file reference, thread colours,
        stitch count, placement, machine, notes, photos.
  \item Workflow: start (records operator and time) $\rightarrow$ complete / request
        approval (records user and time) $\rightarrow$ approve / reject (reason required)
        / request re-swatch (reason required).
  \item Rejected attempts stay visible with reason, photos and timestamps.
  \item Approved attempts are \textbf{immutable}; another sample means a new attempt.
  \item \textbf{Production gate:} if a swatch is required and the latest attempt is not
        approved, every embroidery stage transition (start, progress, complete) is
        rejected by the server with a state error.
  \item Queue states: waiting, swatching, awaiting approval, ready for production,
        in production.
\end{itemize}

\subsection{Department stages and job completion}

\begin{itemize}
  \item Each job has one stage per required department; stages progress independently.
  \item Department users see only their own stages; the job detail page shows all
        statuses for coordination.
  \item An operator may update \textbf{only their own department's stage} (Admin and
        Operations can override).
  \item Stage data: notes, waste quantity, reprint quantity, quantity, progress,
        remaining pieces, start and finish times, completing user.
  \item Completing one stage removes it from that department's Active queue; other
        stages stay open.
  \item Total prints = pieces $\times$ selected print positions.
  \item \textbf{A job only completes when every required stage is complete AND dispatch
        is complete.} Dispatch can never be skipped.
\end{itemize}

\subsection{Readiness traffic lights (computed, never typed in)}

\begin{plainwords}
Every job shows a colour. Green means everything needed to start work is ready: stock
received, screens made, and (for embroidery) the sample approved. White means nothing
that the job actually needs is ready yet. Amber is in between. The colour is calculated
by the server every time something changes -- nobody can set it by hand.
\end{plainwords}

The calculation counts only the gates that apply to the job:

\begin{verbatim}
active gates  = stock? screens? swatch?
green         = all active gates pass
white         = no active gate passes (or job has nothing required yet)
amber         = some pass, some don't
\end{verbatim}

\begin{itemize}
  \item Stock gate: every line received $\ge$ ordered.
  \item Screens gate: made $\ge$ required. A \textbf{blank} required field means
        ``not yet specified'' and therefore \emph{not ready} -- blank is never treated
        as zero.
  \item Swatch gate: not required, or latest attempt approved.
\end{itemize}

\begin{warningbox}[title={Developer note: empty-job trap}]
  The zero-gate case must be checked \emph{before} any gate evaluation:
  \texttt{if (totalRequiredGates === 0) return 'White'}. JavaScript array methods such as
  \texttt{[].every(...)} return \texttt{true} on empty arrays -- relying on them makes a
  job with no products, screens or swatches turn \textbf{Green} the moment it is created.
\end{warningbox}

\subsection{Dispatch and shipments}

\begin{tabularx}{\textwidth}{@{}X X@{}}
\toprule
\textbf{Rule} & \textbf{Detail} \\
\midrule
Four methods with method-specific fields &
  Collection (contact, phone, date, time, confirmation); DPD (consignment, parcels,
  weights, service, instructions, labels); Same Day (courier, driver, time, ETA,
  reference); Fanela Van (driver, route, departure, ETA) \\
Who can edit & Dispatch + Admin/Operations; Office sets planned method/address at entry;
  Director and Packing read-only \\
Booking creates a shipment record & ``Record booking/collection arrangement'' saves a
  shipment; ``record final dispatch'' records departure, user, time \\
Guard rails & Other required stages must be finished first; Collection needs
  confirmation; the Dispatch stage completes only via an explicit \textbf{``Finalise
  dispatch / close shipments''} action -- one dispatched shipment among several must
  never auto-close the stage, and finalise is rejected while any shipment is still open
  (unless explicitly abandoned, with a reason, audited) \\
Part dispatch & While a job is open, one shipment can be dispatched and others follow;
  multiple shipments are kept \\
Cancellation = event, never deletion & Void records a reason against the shipment; the
  booking stays; no carrier-side cancellation unless later integrated \\
\bottomrule
\end{tabularx}

\textbf{Shipment states:}
\begin{center}
\texttt{Draft $\rightarrow$ Booking Arranged $\rightarrow$ Booked $\rightarrow$ Labels
Attached $\rightarrow$ Print Requested $\rightarrow$ Labels Printed $\rightarrow$
Dispatched / Collected}
\end{center}
Any non-final state can move to \texttt{Cancelled/Void} with a reason.

\subsection{DPD -- honest manual mode first}

No live DPD booking in the first release. The manual workflow from the prototype is kept,
but server-stored: book with DPD externally, enter consignment and parcel details, save
the booking \emph{first}, attach returned label files (default limit 25 MB, configurable),
log print requests without ever claiming physical printing succeeded, then confirm
``labels printed'' manually. Booking, tracking and labels survive failures and can be
retried. A future API mode adds real bookings; \textbf{manual mode remains as fallback}.

\subsection{Excel -- useful, but never the truth}

\begin{itemize}
  \item Exports (all jobs, filtered, department, stock/shortage, swatches, shipments,
        audit) are generated on the server and filtered by the caller's permissions.
  \item Historical quantities in reports are snapshots -- \textbf{never summed as stock
        movements}.
  \item Imports are add-only: checked before confirmation, duplicates skipped, imported
        artwork starts as Draft.
  \item Never reconstructed from Excel: historical swatch approvals, shipment records,
        photos, audit history.
\end{itemize}

\subsection{Audit and stock history -- two separate append-only logs}

\begin{enumerate}
  \item \textbf{Stock history} -- receipts, adjustments, corrections with reason, user,
        timestamp. Never edited, never deleted.
  \item \textbf{Operational audit} -- job/header, priority, order lines, artwork,
        screens, swatch, stage, customer master, dispatch, shipment and approval
        changes, with before/after snapshots.
\end{enumerate}
Department users see the entries relevant to their work; sensitive commercial changes
stay restricted.
